\section*{Chapitre \ref{ch:g10:numbers} --- Nombres et ensembles de nombres}

\begin{solution}{exo:g10:numbers:1}
$\frac{15}{3} = 5 \in \N$. \quad $-7 \in \Z$. \quad
$\frac{22}{7} \in \Q$ (ce n'est pas un entier~: $22 = 7 \times 3 + 1$).
\quad $\sqrt 9 = 3 \in \N$. \quad $\sqrt{10} \in \R$ (irrationnel,
puisque $10$ n'est pas le carré d'un rationnel --- admis ici, dans
l'esprit du \cref{thm:g10:numbers:sqrt2}). \quad
$-2.4 = -\frac{24}{10} \in \Q$. \quad $0 \in \N$.
\end{solution}

\begin{solution}{exo:g10:numbers:2}
(a)~$\intco{-2}{5}$~: point plein en $-2$, point creux en $5$.
(b)~$\intoo{3}{+\infty}$~: point creux en $3$, hachures vers la droite.
(c)~$\intoc{-\infty}{-1}$~: hachures depuis la gauche jusqu'à un point
plein en $-1$.
(d)~«~distance de $x$ à $2$ au plus $3$~» signifie
$\abs{x - 2} \leq 3$, c'est-à-dire $x \in \intcc{-1}{5}$~: points
pleins en $-1$ et $5$.
\end{solution}

\begin{solution}{exo:g10:numbers:3}
(a)~Les deux \omterm{def:g10:numbers:interval}{intervalles} se chevauchent entre $0$ et $2$~:
$I \cap J = \intcc{0}{2}$ ($0 \in J$ et $2 \in I$, et chacun
appartient aussi à l'autre ensemble), et
$I \cup J = \intco{-3}{4}$.

(b)~$I \cap J$ est l'ensemble des $x$ tels que $-2 \leq x$ et
$x < 1$~: $\intco{-2}{1}$. La \omterm{def:g10:numbers:interunion}{réunion} couvre tout~:
$I \cup J = \R$.
\end{solution}

\begin{solution}{exo:g10:numbers:4}
$\abs{-6} = 6$~; \quad $\abs{4 - 9} = \abs{-5} = 5$~; \quad
$\abs{-3 - 5} = \abs{-8} = 8$~; \quad
$\sqrt 2 > 1$ donc $\abs{\sqrt2 - 1} = \sqrt2 - 1$~; \quad
$1 - \sqrt2 < 0$ donc $\abs{1 - \sqrt2} = \sqrt2 - 1$ aussi (un
nombre et son opposé ont la même \omterm{def:g10:numbers:abs}{valeur absolue}).
\end{solution}

\begin{solution}{exo:g10:numbers:5}
$\abs{x} = 5$~: distance $5$ à $0$, donc $x = 5$ ou $x = -5$.

$\abs{x - 1} = 3$~: distance $3$ à $1$, donc $x = 4$ ou $x = -2$.

$\abs{x - 4} \leq 1$~: distance au plus $1$ à $4$, donc
$x \in \intcc{3}{5}$.

$\abs{x + 2} < 3$~: distance strictement inférieure à $3$ à $-2$,
donc $x \in \intoo{-5}{1}$.
\end{solution}

\begin{solution}{exo:g10:numbers:6}
Chaque \omterm{def:g10:numbers:interval}{intervalle} est décrit par son centre $a$ (milieu des bornes) et
son rayon $r$ (moitié de la longueur).

$\intcc{2}{8}$~: $a = 5$, $r = 3$, donc $\abs{x - 5} \leq 3$.

$\intoo{-5}{1}$~: $a = -2$, $r = 3$, bornes ouvertes, donc
$\abs{x + 2} < 3$.

$\intcc{-7}{-3}$~: $a = -5$, $r = 2$, donc $\abs{x + 5} \leq 2$.
\end{solution}

\begin{solution}{exo:g10:numbers:7}
Soit $x = 0.272727\dots$. Alors $100x = 27.2727\dots$, et en
soustrayant~:
\[
100x - x = 27.2727\dots - 0.2727\dots = 27,
\]
donc $99x = 27$ et $x = \frac{27}{99} = \frac{3}{11}$, un quotient
d'entiers~: $x$ est rationnel.
\end{solution}

\begin{solution}{exo:g10:numbers:8}
\emph{1.~Vrai~:} ajouter des entiers (nombres entiers positifs ou
négatifs) produit toujours un entier.

\emph{2.~Faux~:} $\frac{1}{2}$ est un quotient des entiers $1$ et $2$
et n'est pas un entier.

\emph{3.~Vrai~:} $\frac pq + \frac{p'}{q'} = \frac{pq' + p'q}{qq'}$
est encore un quotient d'entiers (avec dénominateur non nul).

\emph{4.~Vrai~:} supposons $r$ rationnel, $t$ irrationnel, et
$r + t = s$ rationnel. Alors $t = s - r$ serait une différence de
deux rationnels, donc rationnel (par le~3, appliqué avec $-r$) ---
contradiction. Donc $r + t$ est irrationnel.
\end{solution}

\begin{solution}{exo:g10:numbers:9}
En multipliant $1.414 \leq \sqrt2 \leq 1.415$ par $2 > 0$~:
$2.828 \leq 2\sqrt2 \leq 2.830$.

En ajoutant $3$~: $4.414 \leq \sqrt2 + 3 \leq 4.415$.

En multipliant par $-1 < 0$, les inégalités s'inversent~:
$-1.415 \leq -\sqrt2 \leq -1.414$.
\end{solution}

\begin{solution}{exo:g10:numbers:10}
D'abord le fait auxiliaire. Divisons $p$ par $3$~: le reste est $0$,
$1$ ou $2$, c'est-à-dire $p = 3k$, $p = 3k+1$ ou $p = 3k+2$. En
élevant au carré~:
\[
(3k)^2 = 3(3k^2), \quad
(3k+1)^2 = 3(3k^2 + 2k) + 1, \quad
(3k+2)^2 = 3(3k^2 + 4k + 1) + 1 .
\]
Seul le premier est multiple de $3$~: si $p^2$ est multiple de $3$,
alors $p$ l'est aussi.

Supposons maintenant $\sqrt3 = \frac pq$ irréductible. En élevant au
carré~: $p^2 = 3q^2$, donc $p^2$ est multiple de $3$, donc $p = 3k$.
Alors $9k^2 = 3q^2$, donc $q^2 = 3k^2$ et $q$ est aussi multiple de
$3$ --- mais alors la fraction $\frac pq$ n'était pas irréductible~:
contradiction. Donc $\sqrt3$ est irrationnel.
\end{solution}

\begin{solution}{pb:g10:numbers:1}
\textbf{1.} $-7 \in \Z$~; \ $\frac{13}{4} \in \Q$~; \
$\sqrt{16} = 4 \in \N$~; \ $0.121212\ldots = \frac{12}{99} =
\frac{4}{33} \in \Q$ (devoir du week-end sur les décimales périodiques
du volume précédent)~; \ $\sqrt8 = 2\sqrt2$ est irrationnel~: plus
petit ensemble $\R$~; \ $\pi$~: $\R$.

\textbf{2.} $\frac pq + \frac rs = \frac{ps + qr}{qs}$ et
$\frac pq \times \frac rs = \frac{pr}{qs}$~: des quotients d'entiers
de dénominateur non nul --- donc des rationnels.

\textbf{3.} (a)~Si $r$ est rationnel, $x$ irrationnel et si
$r + x = q$ était rationnel, alors $x = q - r$ serait une différence
de rationnels, donc rationnel (question 2)~: contradiction.
(b)~Si $r \neq 0$ et $r x = q$ rationnel, alors $x = \frac qr$ est
rationnel~: contradiction.

\textbf{4.} $\sqrt2$ et $-\sqrt2$ sont \omterm{ex:g10:numbers:classify}{irrationnels} et de somme $0$~;
$\sqrt2$ et $\sqrt2$ ont pour produit $2$. Des résultats rationnels à
partir d'ingrédients \omterm{ex:g10:numbers:classify}{irrationnels}~: les \omterm{ex:g10:numbers:classify}{irrationnels} ne forment aucun
royaume stable.

\textbf{5.} $\sqrt8 = 2\sqrt2$, donc
$\sqrt2 + \sqrt8 = 3\sqrt2$~: irrationnel (question 3b, multiplicateur
$3$). Et $\sqrt2 \times \sqrt8 = \sqrt{16} = 4 \in \N$~: le produit de
deux \omterm{ex:g10:numbers:classify}{irrationnels} atterrit dans le plus petit royaume de tous.

\textbf{6.} $3.475$, puis $3.471$ (ou $3.4701$, $3.47001$, \dots)~:
ajouter des décimales produit indéfiniment de nouveaux rationnels
entre les deux.

\textbf{7.} Les multiples de $10^{-n}$ jalonnent la droite avec un pas
$10^{-n} < g$. Le premier multiple strictement supérieur à $a$ --- il
existe, puisque les multiples finissent par dépasser $a$ --- se trouve
à un pas au plus au-delà de $a$, donc avant $a + g = b$~: il est
strictement compris entre $a$ et $b$. Or un point de la grille est un
décimal, donc un rationnel~: tout \omterm{def:g10:numbers:interval}{intervalle} de longueur strictement
positive en contient un.

\textbf{8.} Entre $a$ et le rationnel $r_1$ trouvé à la question 7 se
trouve (question 7 de nouveau) un rationnel $r_2$~; entre $a$ et $r_2$
un rationnel $r_3$~; et ainsi de suite~: une infinité, tous distincts,
tous dans l'\omterm{def:g10:numbers:interval}{intervalle} de départ.

\textbf{9.} $r + \frac{\sqrt2}{10^k}$ est la somme d'un rationnel et
d'un irrationnel ($\frac{\sqrt2}{10^k}$ est irrationnel par la
question 3b), donc est irrationnel. Comme
$\frac{\sqrt2}{10^k} < \frac{2}{10^k}$ devient plus petit que toute
borne, pour $k$ assez grand $r + \frac{\sqrt2}{10^k}$ reste avant
$b$~: voilà un irrationnel dans l'\omterm{def:g10:numbers:interval}{intervalle} --- et faire varier $k$
en donne une infinité.

\textbf{10.} Pas de plus petit réel strictement positif~: si $s > 0$
l'était, l'\omterm{def:g10:numbers:interval}{intervalle} $(0, s)$ contiendrait encore un rationnel
(question 7), positif et plus petit que $s$. Pas de nombre juste après
$3$~: tout candidat $c > 3$ laisse l'\omterm{def:g10:numbers:interval}{intervalle} $(3, c)$, qui n'est pas
vide --- le jeu du devoir du week-end sur l'absence de nombre suivant,
devenu théorème sur $\R$.

\textbf{11.} $\abs{x - 5} = 2$~: $x = 3$ ou $x = 7$.
$\abs{x - 5} < 2$~: $x \in \intoo{3}{7}$.
$\abs{x + 1} \geq 3$~: distance à $-1$ au moins égale à $3$, d'où
$x \in \intoc{-\infty}{-4} \cup \intco{2}{+\infty}$.

\textbf{12.} $\abs{d - 80} \leq 0.05$, c'est-à-dire
$d \in \intcc{79.95}{80.05}$. Le piston à $79.97$ est accepté~; celui
à $80.06$ est refusé, à un centième de millimètre près.

\textbf{13.} $(3, -5)$~: $\abs{-2} = 2 \leq 8$. $(-2, -7)$~:
$\abs{-9} = 9 = 2 + 7$, il y a égalité. Même signe~: $\abs{a + b}$ est
la somme des distances, égale à $\abs a + \abs b$. Signes contraires~:
le trajet revient sur ses pas, $\abs{a + b}$ est la \emph{différence}
des distances, strictement plus petite que leur somme (sauf si l'un
des deux est nul). Il y a égalité exactement lorsque $a$ et $b$ sont
de même signe ou que l'un est nul.

\textbf{14.} Les solutions sont les points équidistants de $2$ et de
$-4$~: le milieu, $x = -1$. C'est la médiatrice du devoir du week-end
sur les deux miroirs, réduite à une dimension~: un point unique.

\textbf{15.} $\sqrt{x^2} = \abs{x}$, et non $x$~: pour $x = -3$,
$\sqrt{9} = 3 = \abs{-3} \neq -3$. Alors $x^2 < 9$ se lit
$\sqrt{x^2} < 3$, c'est-à-dire $\abs x < 3$~: $x \in \intoo{-3}{3}$.

\textbf{16.} La mesure affirme seulement que la longueur est dans
l'\omterm{def:g10:numbers:interval}{intervalle} $\intcc{1.41420}{1.41422}$ --- et d'après les questions 7
et 9, \emph{cet \omterm{def:g10:numbers:interval}{intervalle} contient une infinité de rationnels et une
infinité d'\omterm{ex:g10:numbers:classify}{irrationnels}}. Il en va de même pour toute tolérance
future, si petite soit-elle. Aucune mesure ne pourra jamais séparer
les deux familles~; seule une démonstration portant sur la longueur
exacte --- comme celle de la diagonale du carré unité (devoir du
week-end sur l'irrationalité du volume précédent) --- le peut.

\textbf{17.} Elles approchent $\sqrt2$ avec des erreurs inférieures à
$10^{-1}, 10^{-2}, 10^{-3}, \dots$~: les rationnels se pressent contre
$\sqrt2$ à toutes les échelles. Le fait général~: tout \omterm{def:g10:numbers:sets}{nombre réel} est
approché d'aussi près que l'on veut par des rationnels (ses
troncatures décimales) --- la densité encore une fois, vue du côté de
la cible.

\textbf{18.} $a + b \in \intcc{2.4 + 1.1}{2.6 + 1.3} =
\intcc{3.5}{3.9}$~: incertitude $\pm 0.2$ (les erreurs s'ajoutent).
$a \times b \in \intcc{2.4 \times 1.1}{2.6 \times 1.3} =
\intcc{2.64}{3.38}$~: une incertitude d'environ $\pm 0.37$ autour de
$3$ --- pour un produit ce sont les erreurs relatives qui s'ajoutent,
et la précision se dégrade donc plus vite.

\textbf{19.} Écrivons les valeurs exactes $a = 2.5 + e_1$ et
$b = 1.2 + e_2$, avec $\abs{e_1}, \abs{e_2} \leq 0.1$. L'erreur sur la
somme vaut alors $\abs{e_1 + e_2} \leq \abs{e_1} + \abs{e_2} \leq
0.2$~: la règle de l'ingénieur \emph{est} l'inégalité triangulaire de
la question 13.

\textbf{20.} La droite réelle n'a ni trous ni voisins~: après un
nombre, pas de «~suivant~» (question 10). Ses points se répartissent
entre les rationnels --- exactement les décimales finies ou
périodiques (devoir du week-end correspondant du volume précédent) ---
et les \omterm{ex:g10:numbers:classify}{irrationnels}, et les deux familles sont si étroitement
entrelacées que tout \omterm{def:g10:numbers:interval}{intervalle} en contient une infinité de chaque
(questions 8 et 9)~; c'est pourquoi aucune mesure, mais seulement une
démonstration, ne peut les distinguer (question 16). Et la surprise
finale, laissée en promesse~: les \omterm{ex:g10:numbers:classify}{irrationnels} sont plus nombreux que
les rationnels --- non pas en les comptant un à un, mais au sens
précis d'une comparaison entre infinis, théorie construite dans les
volumes universitaires.
\end{solution}
